
Our mechanistic attribution approach tests the hypothesis that LLM advantage decomposes into distinct, quantifiable mechanisms via controlled ablation. We design five sub-hypotheses with falsification criteria: (H-E1) establish lean-auto baseline [10\%, 25\%], (H-M1) NL ablation drops success by $\geq 25$ percentage points, (H-M2) depth filtering drops success by $\geq 5$ percentage points, (H-M3) random Mathlib achieves 18-25\%, (H-C1) tactic budget CV $\leq$ 100\%. Each gate threshold enforces scientific rigor---effects below threshold reject the mechanistic claim.

\subsubsection{3.1 Dataset and Evaluation Protocol}

We use the miniF2F Lean 4 test set (N=244 Olympiad-level problems) with deterministic Lean 4 proof checker validation (binary outcome: valid/invalid). All experiments fix timeout=300s, Mathlib version (Lean 4 compatible), and evaluate on identical problem sets. Success rate is the primary metric (percentage of valid proofs); tactic count serves as a fairness control metric.

The miniF2F Lean 4 subset provides three critical properties for mechanistic ablation: (1) deterministic validation---Lean proof checker eliminates custom extraction pipelines, (2) NL hint availability---problems include informal docstrings/comments amenable to ablation, (3) established difficulty---Olympiad-level ensures non-trivial proof search.

\subsubsection{3.2 H-E1: Baseline Measurement}

\textbf{Objective:} Measure pure automated prover (lean-auto) success rate on miniF2F to establish reference for LLM comparison.

\textbf{Design:} Run lean-auto (Duper backend, deterministic hammer) on all N=244 problems with 300s timeout. Record success rate (primary), tactic count per solved problem, error sources (type elaboration failures, Mathlib API incompatibilities). Duper represents standard automated proving: breadth-first search over Mathlib premises, no LLM, no NL parsing.

\textbf{Predicted outcome:} 15\% [10\%, 25\%] success rate, ~10 tactic evaluations per problem.

\subsubsection{3.3 H-M1: Natural Language Understanding Mechanism}

\textbf{Objective:} Test whether NL hint removal drops LLM success by $\geq 25$ percentage points (validates 60\% contribution claim).

\textbf{Design:} Controlled ablation with paired comparison. For each miniF2F problem, create two variants: (1) NL-intact---original formal statement + docstring/comment hints, (2) NL-ablated---same formal statement, all docstrings/comments stripped. Run LLM on both variants, measure success drop via McNemar paired test. Gate threshold: $\Delta \geq$ 25 percentage points AND p < 0.05.

\textbf{Rationale:} NL ablation isolates linguistic contribution by eliminating natural language signals while preserving formal type signatures. If LLMs rely on NL understanding, success should drop significantly.

\textbf{Predicted outcome:} Baseline (NL-intact) 65\%, ablated 35\%, $\Delta =30$ percentage points [25\%, 35\%].

\subsubsection{3.4 H-M2: Proof Depth Mechanism}

\textbf{Objective:} Test whether proof depth filtering ($\leq 3$ tactics) drops LLM success by $\geq 5$ percentage points (validates 30\% contribution claim).

\textbf{Design:} Post-hoc filtering on LLM-solved problems. Measure success rate on: (1) full dataset---all solved problems regardless of tactic count, (2) shallow subset---only problems solvable with $\leq 3$ tactics. Compute $\Delta$ via McNemar test. Gate threshold: $\Delta \geq$ 5 percentage points.

\textbf{Rationale:} Proof depth filtering tests whether long-range context contributes to LLM advantage. Olympiad problems require multi-step reasoning; if LLMs exploit context maintenance, filtering to shallow proofs should drop success significantly.

\textbf{Predicted outcome:} Full success 65\%, shallow success 50\%, $\Delta =15$ percentage points [5\%, 30\%].

\subsubsection{3.5 H-M3: Corpus Pattern Matching Mechanism}

\textbf{Objective:} Test whether random Mathlib tactic sampling achieves 18-25\% ($\Delta =3-10$ percentage points above lean-auto baseline, validates 10\% contribution).

\textbf{Design:} Random sampler draws tactics from Mathlib corpus with human distribution frequency. For each problem, sample tactics uniformly from distribution until timeout or success. Compare random sampler success vs lean-auto baseline. Gate threshold: success $\in$ [18\%, 25\%].

\textbf{Predicted outcome:} Random Mathlib 20\% [18\%, 25\%], $\Delta =+5$ percentage points [+3\%, +10\%] vs lean-auto.

\subsubsection{3.6 H-C1: Tactic Budget Control}

\textbf{Objective:} Validate tactic evaluation count as fairness metric for LLM vs automated prover comparison (CV $\leq$ 100\%).

\textbf{Design:} Post-hoc analysis on H-E1 (lean-auto) solved problems. Extract tactic count per solved proof, compute coefficient of variation CV = $\sigma / \mu$. If CV << 1.0, tactic count is stable metric. Recommend budget = $\mu + 1\sigma$ for future experiments.

\textbf{Predicted outcome:} Mean tactic count $\mu =10, \sigma =4$, CV=0.40 << 1.0. Recommended budget=15.

\subsubsection{3.7 Statistical Analysis}

\textbf{Paired comparisons (H-M1, H-M2):} McNemar test for correlated binary outcomes. Reports $\chi^2$ statistic, p-value, 95\% confidence interval via bootstrap (10,000 resamples). Effect size: Cohen's h for proportions.

\textbf{Independent comparisons (H-M3 vs H-E1):} Chi-squared test for two independent proportions.

\textbf{Gate decisions:} Mechanistic hypothesis PASS if effect exceeds threshold AND p < 0.05. FAIL if effect below threshold (regardless of significance). Example: H-M2 $\Delta =3.7$\% < 5\% $\rightarrow$ FAIL even though p=0.0027.

All experiments use $\alpha =0.05$ significance level, two-tailed tests. Confidence intervals computed via bias-corrected bootstrap.
